\section{Proofs for structured-argumentation rationality}
\label{sec:annex-p18-structured-rationality}

The source conventions come from the corrected general account and its
corrigendum \cite{ModgilPrakken2013GeneralCorrected,ModgilPrakken2018Corrigendum}.
The return-condition theorem is stated with its actual extra hypothesis;
the countermodel below proves why that hypothesis has not been discarded.

\begin{annexitem}{annex:p18.ratsets}{thm:p18-rationality-strict-sets}{Corrected strict set comparisons}
\begin{proof}
The strict part \(<\) of a preorder is irreflexive and transitive.
An empty set is not strictly below itself by definition.
If a nonempty finite \(\Gamma\) satisfied
\(\Gamma\triangleleft_{\rm E}\Gamma\), its witness \(x\in\Gamma\)
would satisfy \(x<x\), impossible.
If \(\Gamma\triangleleft_{\rm D}\Gamma\), every \(x\in\Gamma\)
would have a strictly larger successor in \(\Gamma\).
Starting at any element and choosing successive larger elements, a finite
set forces a repetition. Transitivity would then imply \(x<x\).
Thus the democratic relation is irreflexive as well; finiteness is
essential in this argument.

Suppose \(\Gamma\triangleleft\Delta\triangleleft\Theta\).
The first two operands must be nonempty. If \(\Theta\) is empty, the
empty-right clause proves \(\Gamma\triangleleft\Theta\).
Otherwise all three are nonempty. In the elitist case choose
\(x\in\Gamma\) below every \(y\in\Delta\), and \(y_0\in\Delta\)
below every \(z\in\Theta\); then \(x<y_0<z\).
In the democratic case, for each \(x\in\Gamma\) choose \(y\in\Delta\)
with \(x<y\), then \(z\in\Theta\) with \(y<z\). In either case
transitivity of \(<\) gives the required comparison.

For the union assertion, the union is nonempty.
If \(\Delta=\varnothing\), any nonempty \(\Gamma_i\) is below it.
If \(\Delta\ne\varnothing\), an elitist witness in the union belongs
to some \(\Gamma_i\) and witnesses its comparison with \(\Delta\).
In the democratic case the universal-existential condition restricts
to every nonempty \(\Gamma_i\). This proves all assertions.
\end{proof}
\end{annexitem}

\begin{annexitem}{annex:p18.ratcounter}{thm:p18-rationality-transposition-boundary}{A transposition-closed countermodel}
\begin{proof}
Use paired literals over \(p,q,r\), ordinary premises \(K_p=\{p,q\}\),
and \(K_n=R_d=\varnothing\). The strict rules consist of both antecedent
permutations of each of the following six rule schemas:
\[
\begin{array}{lll}
 (p,q)\to r,&(p,q)\to\neg r,&(\neg r,q)\to\neg p,\\
 (p,\neg r)\to\neg q,&(r,q)\to\neg p,&(p,r)\to\neg q .
\end{array}
\tag{1}
\]
They are twelve distinct ordered rules. Transposing either of the first
two rules gives two of the other four schemas. Transposing any of the
other four gives a schema already listed, with possibly permuted
antecedents. Thus the complete rule set is closed under transposition.
All conflicts are symmetric, so contrary well-formedness holds.
There are no nullary rules and no axioms, so
\(\operatorname{Cl}_s(K_n)=\varnothing\), which is directly consistent.

The only arguments concluding \(p,q\) are \(P=[p]\), \(Q=[q]\).
For each sign \(\epsilon\in\{+,-\}\), there are exactly two arguments
\(R_{\epsilon,\sigma}\) concluding \(r\) or \(\neg r\), corresponding
to the two orders \(\sigma\) of \(P,Q\).
For each of these four arguments there are two arguments concluding
\(\neg p\), using it and \(Q\) in either order, and two concluding
\(\neg q\), using it and \(P\) in either order. Write them
\(U_{\epsilon,\sigma,\tau}\) and
\(V_{\epsilon,\sigma,\tau}\), with \(\tau\) the final antecedent order.
No rule uses \(\neg p\) or \(\neg q\) as an antecedent.
Consequently the full universe has
\[
 2+4+8+8=22
\]
arguments. \(P,Q\) have their singleton supports; every other argument
has ordinary-premise set \(\{p,q\}\).

Define \(B\prec A\) precisely when
\[
 \operatorname{Prem}(B)=\{p,q\},
 \qquad A\in\{P,Q\}.
 \tag{2}
\]
This is irreflexive, and has no chains of length two, so it is a strict
partial order. If a weak relation is wanted, adjoining identity has
exactly this strict part.

There are no strict firm arguments, since every argument has at least
one ordinary premise. Conditions (i) and (ii) of reasonableness
therefore hold. Profile continuations of a singleton preserve its
ordinary-premise set. Since (2) depends only on that set, they preserve
each comparison and noncomparison, proving (iii).

For (iv), take any nonempty finite set \(\mathcal H\) of arguments.
If it contains an argument \(C_i\) supported by \(\{p,q\}\), no argument
at all is strictly below \(C_i\), by (2). Thus every available omitted-member
continuation satisfies \(C_{-i}\not\prec C_i\).
If it has no such member, it is a subset of \(\{P,Q\}\).
A continuation of the other singleton has that singleton's support,
and is therefore not a full-support argument strictly below the omitted
member. A continuation of the empty family would need empty ordinary
support; no such argument exists. This proves (iv) in every case.
Repeated occurrences of the same tree are not distinct set members.

All arguments are strict, so only undermining attacks occur.
Each \(U_{\epsilon,\sigma,\tau}\) attacks \(P\) and every nonleaf
argument, because all nonleaf arguments contain \(P\).
Each \(V_{\epsilon,\sigma,\tau}\) attacks \(Q\) and every nonleaf
argument. There are no other attacks: the conclusions \(r,\neg r\)
are at strict roots, and no other premise is available.
Every attack compares its full-support source with the targeted
singleton \(P\) or \(Q\). It fails by (2). Hence \(D=\varnothing\).

In an edgeless defeat framework, every argument is defended by every set.
Its only complete extension is therefore the entire universe. That set
is also its grounded, preferred, and stable extension.
It concludes both \(p,\neg p\), both \(q,\neg q\), and both \(r,\neg r\).
This establishes the claimed failure under all stated hypotheses.

The indexed-family condition with repetitions is stronger. For the
four-term family \(P,P,Q,Q\), deleting any one index leaves both
premise types. Choose a full-support argument as each omitted-index
continuation. Each is strictly below its corresponding singleton target
by (2). Thus the order violates that repetition-permitting condition.
The counterexample therefore refutes the finite-set reading, not this
stronger ordering premise.

The example also locates the failed reversal step. Let \(B\) be an
argument for \(\neg p\), and let \(A=P\). Then \(B\prec A\), and
the maximal fallible subarguments of \(B\) are \(P,Q\).
Reversing against \(Q\) as in the source's Proposition 8 would require
a strict argument for \(\neg q\) using only \(P\). No rule in (1)
can apply from the sole premise \(p\), so no such argument exists.
Closure under transposition has not supplied the stronger consequence
contraposition required for that step.
The full-support premises \(\{p,q\}\) are c-inconsistent, so this
countermodel is explicitly an unrestricted framework; it is not asserted
to be a c-restricted framework satisfying c-classicality.
\end{proof}
\end{annexitem}

\begin{annexitem}{annex:p18.ratdefense}{thm:p18-rationality-joint-defense}{From return to attack compatibility}
\begin{proof}
First, no two arguments defended by a defeat-conflict-free set \(E\)
can defeat one another. If \(X D Y\), defense of \(Y\) gives
\(U\in E\) with \(U D X\). Defense of \(X\) gives \(V\in E\)
with \(V D U\), contradicting defeat-conflict-freeness.
This proof also applies if \(X=Y\).

Suppose now that \(X T Y\). The preceding paragraph excludes \(X D Y\).
By (R), choose \(W\) with \(W D X\) such that every defeater of \(W\)
defeats \(X\) or \(Y\). Defense of \(X\) supplies \(U\in E\)
with \(U D W\). Therefore \(U D X\) or \(U D Y\).
Defense of the corresponding argument supplies \(V\in E\) with
\(V D U\), again a contradiction. Interchanging \(X,Y\) excludes the
opposite attack. No admissibility assumption on \(E\), beyond its
defeat-conflict-freeness, was used.
\end{proof}
\end{annexitem}

\begin{annexitem}{annex:p18.ratsupport}{thm:p18-rationality-support-union}{The c-consistent union lemma}
\begin{proof}
Subarguments of a defended argument are defended: any defeater of a
subargument defeats its superargument, with the same target and preference
comparison, so defense of the superargument supplies a counter-defeater.
Also, a tree obtained by grafting defended argument trees into a strict
derivation, with additional axiom leaves if needed, is defended whenever
it belongs to the admitted universe. Its new nodes and added axioms
are not attackable; every defeat lands inside one of the input trees.
These are statements about actual tree grafting, not merely equality
of sets of used rules.

Suppose the finite union \(S=\bigcup_i\operatorname{Prem}(A_i)\) is
c-inconsistent. It is finite because every \(A_i\) is a finite tree.
Choose an inclusion-minimal c-inconsistent subset \(S_0\subseteq S\).
It contains an ordinary premise \(\phi\): otherwise
\(S_0\subseteq K_n\), contradicting axiom consistency.
By c-classicality there is a finite strict derivation of a contradictory
\(\phi^\ast\) from \(S_0\setminus\{\phi\}\).
Use its actual knowledge-base premises as their premise arguments to
obtain an argument \(W\). Its premise set is a subset of the
c-consistent set \(S_0\setminus\{\phi\}\); hence \(W\) is admitted
in the c-restricted universe.
Each premise leaf used in \(W\) is a subargument of some \(A_i\),
and is therefore defended by \(E\). The strict-grafting observation
shows that \(W\) is defended by \(E\).
The ordinary premise argument \([\phi]\) is also a subargument of
some \(A_i\), so it is defended.
But \(W\) undermines \([\phi]\), since its conclusion is contradictory
to \(\phi\). This contradicts
Lemma~\ref{thm:p18-rationality-joint-defense}.
Thus the union is c-consistent.
\end{proof}
\end{annexitem}

\begin{annexitem}{annex:p18.ratgeneral}{thm:p18-rationality-general}{The full return-condition theorem}
\begin{proof}
\emph{Semantic equivalence.}
An attack-conflict-free set is defeat-conflict-free because \(D\subseteq T\).
Conversely, if \(E\) is defeat-admissible, all its members are defended
by the defeat-conflict-free set \(E\). The joint-defense lemma shows
that no two of them attack one another, so \(E\) is attack-admissible.
Thus the two admissible families agree.
Both notions of complete extension additionally require inclusion of every
argument defended through the same relation \(D\), so their complete
families agree. The existing abstract theory then gives the same grounded
and preferred extensions, and the same stable extensions because the
additional condition of defeating every outsider also uses \(D\).
In particular the abstract fundamental lemma and existence of grounded
and preferred extensions apply through this equality of families.

\emph{Subarguments and strict grafts.}
The observations at the beginning of the preceding proof apply equally
to the unrestricted universe. If \(A\in E\) for a complete extension,
each of its subarguments is admitted and defended by \(E\), hence belongs
to \(E\).
If a strict rule has antecedents concluded by \(A_1,\ldots,A_k\in E\),
graft those arguments into that rule. In the unrestricted case the
resulting tree is admitted by definition.
In the c-restricted case its premise set is the union of the input
premise sets, which is c-consistent by the support-union lemma.
Thus it is admitted in this case as well.
Every defeat lands in one of the \(A_i\), so the graft is defended and
belongs to \(E\). A nullary strict argument is unattacked and, by axiom
consistency, has c-consistent empty support. The same conclusion holds.
Closure under one strict-rule application implies closure under all
finite strict derivations by induction on derivation height.
The finite-arity increasing-union construction of strict closure proves
\(\operatorname{Conc}(E)=\operatorname{Cl}_s(\operatorname{Conc}(E))\).

\emph{Producing an attack from contradictory conclusions.}
Let \(X,Y\) be defended by the same defeat-conflict-free set \(E\).
In the restricted case assume their joint support is c-consistent;
this is automatic from the support-union lemma under our hypotheses.
Suppose their conclusions conflict. If the target root is ordinary
or defeasible, there is an immediate attack, contrary to the joint-defense
lemma. If both arguments are strict and firm, both conclusions are in
\(\operatorname{Cl}_s(K_n)\), contrary to axiom consistency.
If the target is strict and firm and the other argument is fallible,
contrary well-formedness makes the conflict symmetric, so interchange
the arguments. We are left with the case in which \(Y\) is fallible
and has a strict top rule. Well-formedness again ensures that
\(\operatorname{Conc}(X)\) and \(\operatorname{Conc}(Y)\) are
contradictories.

For the transposition alternative, select a path from the root of \(Y\)
to a first fallible node \(Y'\): an ordinary leaf or a defeasible
rule node. Such a path exists because \(Y\) is fallible.
Begin with the argument \(W_0=X\), whose conclusion is the paired
contradictory of \(Y\)'s conclusion. At a strict node on the selected
path, transpose its rule at the child on that path.
Supply the other antecedents with the actual sibling subtrees of \(Y\),
and the replaced antecedent with the already constructed \(W_j\).
The resulting strict graft \(W_{j+1}\) concludes the paired contradictory
of the selected child. Antecedent permutations ensure that the rule
uses the correct ordered list. Iterating along the finite path gives
an argument \(W\) concluding the contradictory of
\(\operatorname{Conc}(Y')\). It therefore rebuts or undermines \(Y\).

Every original fallible subtree of \(W\) comes from \(X\) or \(Y\);
only strict nodes were added. Consequently \(W\) is defended by \(E\).
Its premise set is a subset of
\(\operatorname{Prem}(X)\cup\operatorname{Prem}(Y)\), so it is
admitted also in the restricted case. The construction may retain
another occurrence of \(Y'\) in a sibling subtree. It does not claim
the stronger deletion property disproved by the countermodel.
The jointly defended arguments \(W,Y\) now attack, a contradiction.

For the contraposition alternative, cut \(Y\)'s tree at the first
fallible node of each root-to-leaf path. Let \(\mathcal M\) be the
finite set of these actual frontier subtrees; include the remaining
axiom leaves separately. Reading the strict part above the cut gives
\[
 S\vdash_s\operatorname{Conc}(Y),\qquad
 S=\operatorname{Conc}(\mathcal M)\cup\operatorname{Prem}_n(Y).
 \tag{3}
\]
Choose a frontier argument \(Y'\), and put
\(\phi=\operatorname{Conc}(Y')\).
By contraposition, there is a strict derivation of a contradictory
\(\phi^\ast\) of \(\phi\) from
\((S\setminus\{\phi\})\cup\{\operatorname{Conc}(X)\}\).
Replace each premise of this derivation by a corresponding frontier
subtree, axiom leaf, or \(X\). If several subtrees share a conclusion,
choose one; if a premise is unnecessary, no extra subtree is demanded.
This produces a strict graft \(W\) with premises in the joint support
of \(X,Y\), defended by \(E\), and attacking \(Y\) at \(Y'\).
Again it is admitted in either universe and contradicts joint defense.
This argument makes no assertion that every unused resource must remain
in the continuation's component sets.

\emph{Direct and indirect consistency.}
For an admissible \(E\), each pair \(X,Y\in E\) is defended by \(E\).
The preceding argument excludes any conflicting pair of conclusions,
including a self-conflicting conclusion. Thus the conclusions are
directly consistent.
For a complete \(E\), strict closure has already been proved; its
strict closure is the same directly consistent conclusion set and
is therefore indirectly consistent.
All arguments involved in each construction are finite, so no finiteness
assumption on \(E\), \(L\), or the whole argument universe was used.
\end{proof}
\end{annexitem}

\begin{annexitem}{annex:p18.ratexample}{thm:p18-rationality-corrigendum-example}{The complete four-premise example}
\begin{proof}
Only the four positive literals are premises or rule antecedents.
Every rule concludes the negative of the sole positive literal absent
from its antecedents. No rule has a negative antecedent.
Consequently all arguments are the four premise trees and the six
antecedent-order variants \(N_{x,\sigma}\) for each missing literal \(x\);
there are \(4+4\cdot6=28\) in total.
All are strict, with nonempty ordinary support.
The base strict priority relation is empty because the two equivalence
classes are incomparable. No corrected strict set comparison between
nonempty operands holds. Hence either corrected link construction gives
empty strict argument preference and every attack succeeds.
The only attacks are undermining:
\[
 N_{x,\sigma} D P_x,\qquad
 N_{x,\sigma} D N_{y,\tau}\quad (y\ne x).
 \tag{4}
\]
There are no attacks between two variants with the same missing literal.
Every argument has an attacker, so the empty set defends nothing and
is the grounded complete extension.

For each \(x\), the set \(E_x\) is conflict-free by (4).
Its negative arguments defeat the absent premise \(P_x\) and all
negative arguments with a different missing literal. These are exactly
the outsiders, so \(E_x\) is stable.
Any nonempty admissible set must contain a negative argument:
a set of positive premise arguments attacks nothing and therefore cannot
defend any of its members against their negative attackers.
If it contains \(N_{x,\sigma}\), conflict-freeness excludes \(P_x\)
and every \(N_{y,\tau}\) with \(y\ne x\). Thus it is a subset of \(E_x\).
This proves that the four \(E_x\) are exactly the maximal admissible
and stable extensions.
Moreover, the single \(N_{x,\sigma}\) defeats every attacker of all
members of \(E_x\). Any complete extension containing it must therefore
contain all of \(E_x\). Hence the complete list is
\(\varnothing,E_p,E_q,E_r,E_s\).

All attack and defeat edges coincide, so (R) is vacuous.
The strict rules are closed under transposition because transposing at
one antecedent changes the missing positive literal to that antecedent;
the required rule and all permutations are present.
The empty strict closure is empty, and there are no one-way contraries.
For c-classicality, a minimal c-inconsistent set of paired literals is
either a contradictory pair \(\{x,\neg x\}\) or
\(\{p,q,r,s\}\). To verify this list, observe that strict inference can
only add negative literals; if a set has no contradictory pair and does
not contain all four positives, a newly derived \(\neg x\) has its
positive \(x\) absent. It therefore remains c-consistent.
A contradictory pair has the c-classical property by its opposite
singleton leaf. In the four-positive set, removing \(x\) leaves the
antecedents of the rule deriving \(\neg x\). This proves c-classicality.
Every generated rule argument has only three positive premises, which
are c-consistent. Thus restricting to c-consistent arguments removes
none of the 28 arguments.

Finally apply the tutorial's weak elitist relation instead.
The support \(\{p,q,r,s\}\setminus\{x\}\) of \(N_{x,\sigma}\)
contains the other member of \(x\)'s equivalence class, so it weakly
compares below \(\{x\}\). The reverse comparison fails because the
three-element support also contains members of the incomparable class.
The asymmetric part therefore puts every \(N_{x,\sigma}\) strictly
below its target premise \(P_x\). Every undermining attack in (4)
fails at precisely that target, so \(D=\varnothing\).
All 28 arguments then form the unique defeat-complete, grounded,
preferred, and stable extension. Their conclusions contain all four
positive and negative pairs. This is the corrigendum's distinction
between the two priority constructions, calculated with every tree
and every edge present.
\end{proof}
\end{annexitem}
